\label{sec:intervals-as-subsets}

\begin{topicbox}{Intervals as the Basic Pieces}
\begin{exposition}
The interval shapes were introduced in \S\ref{sec:intervals-on-a-number-line}. Here they are recast as subsets of the metric line and supplemented with boundedness, which is the size condition compactness will need.
\end{exposition}
\end{topicbox}
\begin{definitionbox}{Definition (Bounded Subset of the Real Line)}
\begin{definition}[Bounded Subset of the Real Line]
\label{def:bounded-subset-of-r}
A set $E\subseteq\mathbb{R}$ is \emph{bounded} exactly when there exists $M\ge 0$ such that $|x|\le M$ for all $x\in E$; equivalently, $E$ is contained in some closed interval $[-M,M]$.
\end{definition}
\end{definitionbox}
\begin{remark*}[Standard quantified statement]
\[
E \text{ is bounded} \Longleftrightarrow \exists M\ge 0\;\forall x\in E\;(|x|\le M).
\]
\end{remark*}
\begin{remark*}[Negated quantified statement]
\[
E \text{ is unbounded} \Longleftrightarrow \forall M\ge 0\;\exists x\in E\;(|x|>M).
\]
\end{remark*}
\begin{remark*}[Interpretation]
Boundedness says $E$ fits inside a single closed interval --- it does not run off to infinity. Together with closedness it is one half of the Heine--Borel characterization of compact subsets of $\mathbb{R}$.
\end{remark*}
\begin{dependencies}
\begin{itemize}
  \item \hyperref[def:closed-interval]{Closed Interval}
  \item \hyperref[thm:triangle-inequality]{Triangle Inequality}
\end{itemize}
\end{dependencies}

\subsection*{Set Operations on Intervals}
\label{subsec:set-operations-on-intervals}

\begin{topicbox}{Intervals as Sets}
\begin{exposition}
Once intervals are regarded as subsets of $\mathbb{R}$, the usual set
operations apply to them. This is the first place where intervals behave like
the building blocks of topology rather than merely as pieces of notation on
the number line. Intersections often remain intervals, while unions,
complements, and differences may split into several interval pieces.
\end{exposition}
\end{topicbox}

\begin{definitionbox}{Definition (Union of Intervals)}
\begin{definition}[Union of Intervals]
\label{def:union-of-intervals}
Let $I,J\subseteq\mathbb{R}$ be intervals. Their \emph{union} is the set
\[
I\cup J=\{x\in\mathbb{R}:x\in I\text{ or }x\in J\}.
\]
\end{definition}
\end{definitionbox}

\begin{dependencies}
\begin{itemize}
  \item \hyperref[def:open-interval]{Open Interval}
  \item \hyperref[def:closed-interval]{Closed Interval}
  \item \hyperref[def:union]{Union}
\end{itemize}
\end{dependencies}

\begin{definitionbox}{Definition (Intersection of Intervals)}
\begin{definition}[Intersection of Intervals]
\label{def:intersection-of-intervals}
Let $I,J\subseteq\mathbb{R}$ be intervals. Their \emph{intersection} is the set
\[
I\cap J=\{x\in\mathbb{R}:x\in I\text{ and }x\in J\}.
\]
\end{definition}
\end{definitionbox}

\begin{dependencies}
\begin{itemize}
  \item \hyperref[def:open-interval]{Open Interval}
  \item \hyperref[def:closed-interval]{Closed Interval}
  \item \hyperref[def:intersection]{Intersection}
\end{itemize}
\end{dependencies}

\begin{definitionbox}{Definition (Complement of an Interval)}
\begin{definition}[Complement of an Interval]
\label{def:complement-of-interval}
Let $I\subseteq\mathbb{R}$ be an interval. Its \emph{complement} in
$\mathbb{R}$ is the set
\[
I^{c}=\mathbb{R}\setminus I
=\{x\in\mathbb{R}:x\notin I\}.
\]
\end{definition}
\end{definitionbox}

\begin{dependencies}
\begin{itemize}
  \item \hyperref[def:open-interval]{Open Interval}
  \item \hyperref[def:closed-interval]{Closed Interval}
  \item \hyperref[def:complement]{Complement}
\end{itemize}
\end{dependencies}

\begin{definitionbox}{Definition (Difference of Intervals)}
\begin{definition}[Difference of Intervals]
\label{def:difference-of-intervals}
Let $I,J\subseteq\mathbb{R}$ be intervals. The \emph{difference} of $I$ by
$J$ is the set
\[
I\setminus J=\{x\in\mathbb{R}:x\in I\text{ and }x\notin J\}
=I\cap J^{c}.
\]
\end{definition}
\end{definitionbox}

\begin{remark*}[Predicate readings]
For any $x\in\mathbb{R}$,
\[
x\in I\cup J \Longleftrightarrow (x\in I)\lor(x\in J),
\]
\[
x\in I\cap J \Longleftrightarrow (x\in I)\land(x\in J),
\]
\[
x\in I^c \Longleftrightarrow x\notin I,
\qquad
x\in I\setminus J \Longleftrightarrow (x\in I)\land(x\notin J).
\]
These are set-theoretic operations; no arithmetic operation is being performed
on the endpoints.
\end{remark*}

\begin{dependencies}
\begin{itemize}
  \item \hyperref[def:open-interval]{Open Interval}
  \item \hyperref[def:closed-interval]{Closed Interval}
  \item \hyperref[def:left-half-open-interval]{Left-Half-Open Interval}
  \item \hyperref[def:right-half-open-interval]{Right-Half-Open Interval}
  \item \hyperref[def:union-of-intervals]{Union of Intervals}
  \item \hyperref[def:intersection-of-intervals]{Intersection of Intervals}
  \item \hyperref[def:complement-of-interval]{Complement of an Interval}
\end{itemize}
\end{dependencies}

\begin{remark*}[Intersections of intervals]
The intersection of two intervals is either empty or another interval. For
example, if $a\leq b$ and $c\leq d$, then
\[
[a,b]\cap[c,d]
=
\begin{cases}
[\max\{a,c\},\min\{b,d\}], & \max\{a,c\}\leq \min\{b,d\},\\
\varnothing, & \max\{a,c\}>\min\{b,d\}.
\end{cases}
\]
Endpoint inclusion follows the endpoint rules of the original intervals.
\end{remark*}

\begin{remark*}[Unions of intervals]
The union of two intervals need not be an interval:
\[
(0,1)\cup(2,3)
\]
has a gap. If two intervals overlap or touch without a missing endpoint, their
union is again an interval, for instance
\[
(0,2]\cup[2,5)=(0,5).
\]
The missing or included endpoints determine whether the joined set is truly
one interval.
\end{remark*}

\begin{remark*}[Complements of intervals]
The complement of a bounded interval usually splits into two rays:
\[
(a,b)^c=(-\infty,a]\cup[b,\infty),
\qquad
[a,b]^c=(-\infty,a)\cup(b,\infty).
\]
Complements reverse endpoint inclusion: an endpoint excluded from the interval
is included in the complement, and an endpoint included in the interval is
excluded from the complement.
\end{remark*}

\begin{remark*}[Differences of intervals]
A difference is an intersection with a complement, so it can also split:
\[
(0,5)\setminus[2,3]=(0,2)\cup(3,5).
\]
If the removed interval only cuts off one side, the result may remain a single
interval:
\[
(0,5)\setminus(3,\infty)=(0,3].
\]
The bracket at $3$ appears because $3\notin(3,\infty)$, so it is not removed.
\end{remark*}

\begin{remark*}[Topological motivation]
Open sets in $\mathbb{R}$ are built from unions of open intervals. Closed sets
are understood through complements of open sets. Boundaries arise when a point
is not cleanly inside a set or its complement. Thus interval set operations are
the grammar behind the metric topology of the real line.
\end{remark*}
